\documentclass[11pt]{article}
\usepackage[margin=1in]{geometry}
\usepackage{amsmath,amssymb}

\newcommand{\gbar}{\bar{g}}

\begin{document}
\begin{center}
{\large\bfseries The Laplace approximation of $\log p(r)$}
\end{center}
\medskip

\section{The quantity}
At a query image the GP gives a Gaussian posterior on the log firing rate $g$:
\[
  g \sim \mathcal{N}(\mu_g,\sigma_g^2),\qquad
  \mu_g = A\,\bar\lambda + \lambda_0,\quad \sigma_g^2 = A^2\,\mathrm{var}(\lambda).
\]
Given $g$ the spike count is Poisson with rate $e^{g}$. The predictive
probability of $r$ spikes marginalises over $g$:
\[
  p(r) = \int_{-\infty}^{\infty}
         \underbrace{\frac{e^{rg}\,e^{-e^{g}}}{r!}}_{\mathrm{Poisson}(r;\,e^{g})}\,
         \underbrace{\frac{1}{\sqrt{2\pi\sigma_g^2}}\,
           e^{-(g-\mu_g)^2/2\sigma_g^2}}_{\mathcal{N}(g;\,\mu_g,\sigma_g^2)}
         \,dg .
\]
This integral has no closed form.

\section{The integrand and its peak}
Write the integrand as $e^{\psi(g)}$ up to constants, with
\[
  \psi(g) = r\,g - e^{g} - \frac{(g-\mu_g)^2}{2\sigma_g^2}.
\]
$\psi$ is strictly concave, so the integrand has a single peak at $\gbar$,
where $\psi'(\gbar)=0$:
\[
  r = e^{\gbar} + \frac{\gbar-\mu_g}{\sigma_g^2}.
\]
This solves in closed form through the principal Lambert-$W$ function:
\[
  \gbar(r) = r\sigma_g^2 + \mu_g - W_0\!\big(\sigma_g^2\,e^{\,r\sigma_g^2+\mu_g}\big).
\]
(The code evaluates $W_0$ in log-space, so the exponential is never formed.)

\section{The approximation}
Laplace replaces $\psi$ by its second-order Taylor expansion at the peak,
\[
  \psi(g)\approx \psi(\gbar) + \tfrac12\,\psi''(\gbar)\,(g-\gbar)^2,
\]
i.e.\ a Gaussian, whose integral is analytic. Carrying the constants through:
\[
  \boxed{\;\log p(r) \approx
    r\,\gbar - e^{\gbar}
    - \frac{(\gbar-\mu_g)^2}{2\sigma_g^2}
    - \tfrac12\log\!\big(1+\sigma_g^2 e^{\gbar}\big)
    - \log r!\;}
\]
Every term is evaluated \emph{exactly} at $\gbar$. This expression is the
output the algorithm uses, and the object whose accuracy we test.

\section{Where the error comes from}
The error is exactly what the quadratic expansion drops: the third and higher
derivatives of $\psi$ at the peak (its asymmetry). Split $\psi$ into its two
factors and differentiate:
\[
  \psi(g) = \underbrace{r g - e^{g}}_{\text{Poisson}}
          \;\underbrace{-\,(g-\mu_g)^2/2\sigma_g^2}_{\text{Normal}},
\qquad
\begin{aligned}
  \psi''(g)   &= -e^{g} - 1/\sigma_g^2,\\
  \psi'''(g)  &= -e^{g},\\
  \psi''''(g) &= -e^{g}.
\end{aligned}
\]
The Normal factor is exactly quadratic: it adds the constant $-1/\sigma_g^2$ to
the curvature and \emph{zero} to every higher derivative. The Gaussian fit
reproduces it exactly, so it contributes nothing to the error. All of the error
comes from the $-e^{g}$ (Poisson) term, the only non-quadratic part of $\psi$.

\section{Why $\sigma_g^2$ controls the error}
$\sigma_g^2$ enters only through the curvature at the peak,
\[
  |\psi''(\gbar)| = e^{\gbar} + \frac{1}{\sigma_g^2},
\]
which sets the width of the fitted Gaussian, $1/\sqrt{|\psi''(\gbar)|}$.

Because the integrand's leading factor is the Gaussian
$e^{\frac12\psi''(\gbar)(g-\gbar)^2}$, its weight lies within this width of the
peak. The dropped cubic remainder $\tfrac16\psi'''(\gbar)(g-\gbar)^3+\dots$ --- a
fixed coefficient $\psi'''(\gbar)=-e^{\gbar}$ times $(g-\gbar)^3$ --- therefore
matters only at $(g-\gbar)$ of that size, where its magnitude is
\[
  \tfrac16\,\frac{|\psi'''(\gbar)|}{|\psi''(\gbar)|^{3/2}}
  = \tfrac16\,\frac{e^{\gbar}}{\big(e^{\gbar}+1/\sigma_g^2\big)^{3/2}} .
\]
So the asymmetry is measured against the curvature, and $\sigma_g^2$ enters the
error only through this width, leaving $\psi'''$ untouched.
\begin{itemize}
\item Small $\sigma_g^2$: $1/\sigma_g^2$ dominates, the peak is narrow, and over
that width $e^{g}$ barely curves — its asymmetry is invisible and the
approximation is near-exact.
\item Large $\sigma_g^2$: the confinement weakens, the peak widens, and the
asymmetry of $e^{g}$ is exposed — the error grows.
\end{itemize}
So $\sigma_g^2$ is the knob and the Poisson $e^{g}$ is the source: the prior
does not cause the error, it suppresses it by pinning the peak.

\medskip
\noindent\textit{The difference $\log p_{\mathrm{true}}(r)-\log p_{\mathrm{Laplace}}(r)$
is what \texttt{laplace\_pointwise\_error.py} measures against an
adaptive-quadrature reference.}

\end{document}
